%% Bilan des forces sur un livre posé sur une table (leçon pc-forces).
%% Deux forces seulement : le poids P (vers le bas) et la réaction R de la table
%% (vers le haut). Elles sont tracées de MÊME LONGUEUR et sur la même verticale :
%% c'est la traduction graphique de l'équilibre R = P.
\documentclass[tikz,border=4pt]{standalone}
\usepackage{liaisons}
\begin{document}
\begin{tikzpicture}[x=1cm,y=1cm,
  vect/.style={-{Stealth[length=6pt]}, line width=1.2pt}]

%% ------------------- la table -----------------------------------------
\filldraw[fill=black!15, draw=black!55, line width=0.7pt]
      (-3,-0.35) rectangle (3,0);
\node[font=\small, anchor=north west] at (-3,-0.42) {Table};

%% ------------------- le livre -----------------------------------------
\filldraw[fill=solideD!30, draw=solideD, line width=1pt]
      (-0.8,0) rectangle (0.8,0.7);
\node[font=\small, text=solideD!75!black, anchor=west] at (0.1,0.35) {Livre};

%% ------------------- centre de masse ----------------------------------
\coordinate (G) at (0,0.35);

%% ------------------- les deux forces ----------------------------------
% réaction de la table : verticale, vers le haut, longueur 1,6 cm
\draw[vect, solideE] (G) -- ++(0,1.6)
      node[anchor=west, inner sep=3pt, font=\small, text=solideE] {$\vec{R}$};
% poids : verticale, vers le bas, MÊME longueur 1,6 cm
\draw[vect, black] (G) -- ++(0,-1.6)
      node[anchor=west, inner sep=3pt, font=\small] {$\vec{P}$};
\fill[black] (G) circle (1.6pt);
\node[font=\small, anchor=east, inner sep=4pt] at (G) {$G$};

%% ------------------- cartouche ----------------------------------------
\node[font=\footnotesize, anchor=north] at (0,-2.05)
     {Système : le livre \quad---\quad Référentiel : terrestre};
\end{tikzpicture}
\end{document}
